%=== CHAPTER ONE ===
\chapter{Introduction}
\label{ch:introduction}
\begin{spacing}{1.5}
\setlength{\parskip}{0.22in}

\section{Background}

Quadrotors are increasingly used for inspection, transport and delivery because they can hover and manoeuvre in confined three-dimensional environments. Their agility, however, is obtained from a light airframe with limited thrust and torque margins. A package pickup can therefore change the vehicle dynamics on the same time scale as the control loop. The total mass changes the thrust-to-acceleration map, an offset load produces a persistent moment about the composite centre of mass (CoM), and the payload lever arm changes the rotational inertia.

Model predictive control (MPC) is attractive for this problem because it predicts future motion while handling actuator limits and coupled dynamics explicitly \cite{rawlings2017mpc,mayne2014mpc}. For a quadrotor, the relevant equations are nonlinear and the attitude evolves on the unit-quaternion manifold; nonlinear MPC (NMPC) therefore avoids the loss of fidelity caused by repeated linearisation over aggressive motion \cite{grune2017nmpc,kamel2017linear}. Its performance nevertheless depends on the model supplied to the prediction. If the controller continues to use the empty-airframe mass after grasping a parcel, even a numerically converged optimisation solves the wrong problem.

This dissertation couples moving horizon estimation (MHE) with NMPC. MHE estimates states or parameters by optimising over a recent window of measurements while enforcing the system dynamics and numerical bounds. No task-specific parcel-mass value is supplied to the adopted MHE. Instead, physical rotor thrust establishes an unloaded estimate before NMPC takeover, a pickup produces a new composite-mass estimate, and their difference provides the parcel-mass prediction. A payload step is problematic because a window temporarily contains data generated by two different plants while the estimated mass is constrained to be constant within that window. The implemented event rule therefore deweights pre-event stages once a persistent rotor-thrust change confirms the discontinuity. At drop, the obsolete payload geometry is released without assigning the empty-airframe mass to the estimator; subsequent measurements must drive the mass state back to its unloaded value. The resulting mass and payload-geometry information is fed back into both the prediction model and the hover-input baseline of the NMPC. This estimator--controller interconnection is referred to throughout as the \emph{MHE--NMPC architecture}.

\section{Problem Statement}

Consider a quadrotor that follows a three-dimensional figure-eight path, grasps an eccentric payload, transports it and releases it at a selected phase of the trajectory. The controller has no access to the true payload parameters during the adaptive run. Only the normal onboard or simulated flight signals are available: odometry, attitude, body rate, actuator/motor-speed information and the commands issued to the flight stack.

Three model errors appear at attachment:
\begin{enumerate}
    \item the composite mass changes abruptly, shifting the hover thrust and the translational acceleration gain;
    \item the payload lever arm increases roll and pitch inertia, reducing inner-loop bandwidth;
    \item the horizontal CoM offset creates a non-zero trim moment and consumes part of the available torque authority.
\end{enumerate}
At release, all three changes occur in the opposite direction. A useful adaptive controller must identify the change rapidly, maintain physical constraints during the transition and avoid using a measurement that is already a function of the current estimate.

The central research question is therefore:
\begin{quote}
Can a quadrotor predict an unknown parcel mass from its existing flight signals, without receiving a task-specific parcel-mass value, and make the estimate return to the unloaded regime after release while preserving constrained NMPC trajectory tracking?
\end{quote}

\section{Research Objectives}

The dissertation has six objectives:
\begin{enumerate}
    \item derive a consistent 13-state nonlinear quadrotor model with explicit payload mass, inertia and CoM entry points;
    \item formulate an MHE in which composite mass is a fourteenth state with no empty-airframe or parcel-mass prior in its arrival weight, initialisation or lower bound;
    \item predict parcel mass from the estimated composite-mass increment, detect pickup/drop from physical rotor thrust, and recover the unloaded mass after drop without resetting the estimator state;
    \item formulate an NMPC whose dynamics and equilibrium-input penalty are updated from the live estimates;
    \item implement the estimator and controller as ROS~2 nodes using acados, PX4 SITL, Gazebo and MAVROS;
    \item evaluate estimation accuracy, tracking performance, constraint margin and computation time under mass steps, eccentric pickup/release events and a three-dimensional figure-eight trajectory, including prior-free, fixed-weight and L1-adaptive baselines.
\end{enumerate}

\section{Scope}

The work focuses on simulation and software-in-the-loop validation of a rigidly attached parcel. Airframe inertia, rotor geometry and thrust coefficient are treated as calibrated plant quantities. The calibrated airframe mass may appear in the rigid-body model and in offline scoring, but it is not used to seed, anchor or lower-bound the MHE mass state. No parcel-specific mass value or joint-state truth is supplied to the estimator. Composite mass is estimated online and parcel mass is reported as its increase over the unloaded model mass. The coupled MHE can generate the roll/pitch inertia increment internally from the mass decision and attachment geometry; revision \texttt{f2264b3} also provides a bounded NMPC-side tracker that maps the published mass prediction into $\Delta J$. The latter retains a \SI{0.15}{kg} lift-safety floor and a \SI{0.5}{kg} platform envelope for initial inner-loop gain scaling. These are non-task-specific safety settings, but they are mass-valued information and are therefore exposed as ablation variables rather than hidden under the phrase ``no prior''. The horizontal CoM offset can be obtained from rotor-speed torque balance. Aerodynamic drag is represented by a linear term.

The scope does not include deformable or sloshing loads, cable-suspended payload dynamics, motor failure, wind-tunnel identification or formal closed-loop stability proof for the complete estimator--controller interconnection. These limitations are revisited in Chapter~\ref{ch:conclusion}.

\section{Contributions}

The contributions are:
\begin{enumerate}
    \item a no-initial-mass MHE that removes the empty-airframe mass from the estimator lower bound and arrival cost, seeds from measured physical thrust or a numerical-box midpoint, and publishes a parcel-mass prediction rather than consuming a task-specific parcel value;
    \item a non-circular physical-input reconstruction from rotor speeds, avoiding the loss of observability that occurs when the controller's own mass-dependent intended thrust is fed back to the estimator;
    \item an event-triggered MHE weight schedule that removes stale pre-event evidence without discarding the state-history warm start, together with a signed thrust-step detector that distinguishes pickup from drop;
    \item a drop-regression mechanism that releases obsolete geometry but never hard-resets mass, requiring the estimator to return to the unloaded solution from measurements;
    \item a bounded online inertia map that propagates the predicted parcel mass into $\Delta J$, with explicit floor, cap and ratchet ablations and a corrected inner-loop gain-rescaling baseline;
    \item a parameter-consistent NMPC cost in which the hover thrust and trim torque move with the estimated payload;
    \item an end-to-end ROS~2/PX4/Gazebo test architecture with independent mass-step and magnetic-gripper scenario lines and an L1-adaptive comparison;
    \item a reproducible evaluation matrix that separates cold-start observability, parcel-mass prediction, drop regression, closed-loop benefit and event-signal dependence.
\end{enumerate}

\section{Dissertation Organisation}

Chapter~\ref{ch:literature} reviews MPC, NMPC, online estimation and payload-adaptive quadrotor control. Chapter~\ref{ch:model} derives the model and the payload parameter structure. Chapter~\ref{ch:framework} presents the coupled MHE--NMPC architecture, event-triggered estimator and numerical formulation. Chapter~\ref{ch:evaluation} describes the ROS~2/PX4 implementation, paired experiment matrix, metrics and current SITL evidence. Chapter~\ref{ch:conclusion} summarises the findings and identifies the steps required for evidence freeze and hardware deployment.

\end{spacing}
\newpage
